\begin{perchbox}
\textbf{The proofs trust} Lean's kernel (\fact{lean-rule-core};
\fact{lean-perch-canon} for the canonical-encoding package) and its axioms,
\fact{standard-axioms}; Charon \fact{charon-rev}, Aeneas \fact{aeneas-rev} and
Aeneas's Lean library, which models Rust's primitives; the MIR of rustc
\fact{rustc-mir-nightly}; the specification, that is the denotations and the
document meaning \code{Doc.allows}, written from Perch's earlier account at
\fact{doc-allows-source}; the host models, tested against
\fact{host-model-tested-versions}; and, for
Theorem~\ref{stmt:canonical-identity}, that its Lean emitter matches Perch's
\code{canon.rs}, which golden fixtures and fuzzing test. \textbf{The deployed
account also trusts} rustc \fact{rustc-perch} and its WebAssembly back end; the
document compiler; the tested account shell, verifiers and recovery controller;
and soroban-env-host.
\end{perchbox}

\paragraph{Cryptographic assumptions.} Lemma~\ref{lem:binding} assumes an
EUF-CMA-secure scheme and a collision-resistant SHA-256, and
Theorem~\ref{stmt:canonical-identity} gives hash identity only up to SHA-256
collisions. Ed25519~\cite{rfc8032} is checked with strict verification; P-256 is
ECDSA with low $S$ enforced; the WebAuthn verifier~\cite{webauthn2} checks
neither the relying-party ID nor the origin; \fact{mldsa-verifier} is pure
FIPS~204~\cite{fips204} with a fixed context string. The compiler requires keys
in verifier-canonical form, but thresholds count configured identities: a
verifier-address/key pair for an external signer, or an address for a delegated
one. Distinct identities may share a physical key or controlling party;
independent custody is an authoring assumption.
The deployed Ed25519 and WebAuthn verifiers trap on a bad signature
(\fact{verifier-trap-error} for Ed25519). The zero-knowledge recovery factor
(\perchdoc{vision}) proves knowledge of a secret behind a public Poseidon2
commitment, using an UltraHonk proof with KZG commitments over BN254 and
Fiat--Shamir with Keccak; Perch's documents do not name the setup ceremony. The
secret needs at least \fact{zk-secret-entropy-bits} bits of entropy, and Perch's
zero-knowledge verifier is unaudited. Ed25519, P-256 and BN254 are not
post-quantum, so an \fact{mldsa-verifier} signer resists a quantum adversary only
if every path that can change the account's rules does. A contract ID commits to
its deployer and salt, not its code, so calling fixed infrastructure code assumes
that the registry deploys each pinned Wasm hash at the salt derived from it and
that those contracts cannot replace their code. rule-core's claims ledger lists
its \fact{claims-assumed} assumed and \fact{claims-unproven} unproven claims
(Appendix~\ref{sec:reproduce}).
